\chapter{Υλοποίηση σε \tl{GPU}}
\label{ch:gpu}

\section{Σχεδιαστικές επιλογές}
\label{sec:gpu-design}

\subsection{Αντιστοίχιση σε πλέγμα και ομάδες νημάτων}
\label{sec:gpu-mapping}

\section{Διτονική ταξινόμηση}
\label{sec:gpu-bitonic}

Το χρονοδιάγραμμα των επιπέδων είναι το ίδιο με εκείνο της
ενότητας~\ref{sec:cpu-bitonic}, καθώς κατασκευάζεται από κοινό κώδικα. Ό,τι
αλλάζει είναι ο τρόπος εκτέλεσής του, και το πρόβλημα που καλείται να λύσει η
υλοποίηση παραμένει το ίδιο: το πλήθος των επιπέδων είναι μεγάλο και η αφελής
αντιστοίχιση ενός επιπέδου σε μία εκκίνηση πυρήνα θα σήμαινε δεκάδες διαδοχικά
περάσματα πάνω από την καθολική μνήμη.

Στη μονάδα γραφικών προστίθεται ένα δεύτερο κόστος που δεν υπάρχει στον
επεξεργαστή. Κάθε εκκίνηση πυρήνα έχει σταθερή επιβάρυνση, και μια σάρωση με
εκατοντάδες εκκινήσεις πληρώνει την επιβάρυνση αυτή εκατοντάδες φορές. Για τον
λόγο αυτό ολόκληρο το δίκτυο δεν εκκινείται πυρήνα προς πυρήνα, αλλά
καταγράφεται μία φορά ως γράφος εκτέλεσης και στη συνέχεια αναπαράγεται. Η
δρομολόγηση γίνεται έτσι μία φορά αντί για κάθε επίπεδο.

Οι πυρήνες που συνθέτουν τον γράφο είναι τριών ειδών, και η επιλογή μεταξύ τους
καθορίζεται από την απόσταση $j$ του επιπέδου, όπως ακριβώς και στον
επεξεργαστή. Όταν η απόσταση είναι μικρή, μια ολόκληρη ακολουθία επιπέδων
εκτελείται μέσα στην κοινή μνήμη μιας ομάδας νημάτων, όπως περιγράφεται στην
ενότητα~\ref{sec:gpu-bitonic-tiles}. Όταν είναι πολύ μεγάλη ώστε να μη χωρά σε
πλακίδιο, τα δεδομένα παραμένουν στην καθολική μνήμη και ένας πυρήνας εκτελεί
έως πέντε διαδοχικά επίπεδα κρατώντας τα ενδιάμεσα αποτελέσματα σε καταχωρητές.
Τα επίπεδα περικοπής, τέλος, έχουν δικό τους πυρήνα, ο οποίος διαβάζει δύο
στοιχεία και γράφει ένα σε δεύτερο πίνακα, με τους δύο πίνακες να
εναλλάσσονται όπως και στον επεξεργαστή.

\subsection{Οργάνωση σε πλακίδια και κοινή μνήμη}
\label{sec:gpu-bitonic-tiles}

Η κοινή μνήμη μιας ομάδας νημάτων είναι ρητά διαχειριζόμενη, σε αντίθεση με την
κρυφή μνήμη του επεξεργαστή που λειτουργεί αυτόματα. Η υλοποίηση φορτώνει ένα
πλακίδιο δεδομένων σε αυτήν, εκτελεί όσα επίπεδα είναι δυνατόν εξ ολοκλήρου
μέσα της, και το γράφει πίσω μία φορά. Όσα επίπεδα συγχωνευθούν με τον τρόπο
αυτό δεν κοστίζουν καμία επιπλέον κίνηση προς την καθολική μνήμη.

Το μέγεθος του πλακιδίου προκύπτει από δύο περιορισμούς. Αφενός δεν μπορεί να
ξεπεράσει τα 48 \tl{KB} που διατίθενται ανά ομάδα νημάτων χωρίς ειδική
ρύθμιση. Αφετέρου επιλέγεται μικρότερο όταν το σύνολο εργασίας χωρά στην κρυφή
μνήμη δεύτερου επιπέδου της συσκευής, οπότε ο περιορισμός δεν είναι πλέον η
κίνηση προς την κύρια μνήμη αλλά το πλήθος των ομάδων που μπορούν να είναι
ταυτόχρονα ενεργές. Για ακεραίους των 32 \tl{bit} το πλακίδιο καταλήγει στα 8192
στοιχεία.

Το όριο αυτό έχει συνέπεια που καθορίζει τη συμπεριφορά ολόκληρης της
υλοποίησης. Ένα επίπεδο μπορεί να συγχωνευθεί μόνο αν η απόσταση των στοιχείων
που συγκρίνει χωρά στο πλακίδιο, δηλαδή αν είναι μικρότερη από το ήμισυ του
μεγέθους του. Επειδή οι αποστάσεις της πρώτης φάσης φτάνουν έως το ήμισυ του
μεγέθους των ταξινομημένων μπλοκ, το κρίσιμο μέγεθος είναι το $k'$ της
ενότητας~\ref{sec:cpu-bitonic}: όσο τα μπλοκ χωρούν στο πλακίδιο, σχεδόν όλα τα
επίπεδα συγχωνεύονται.

Για τις μικρές και μεσαίες τιμές του \tl{k} αυτό ισχύει, και η κίνηση δεδομένων
προκύπτει ουσιαστικά ίδια με εκείνη του επεξεργαστή: περίπου επταπλάσια της
αφελούς εκτίμησης, με τις δύο αρχιτεκτονικές να συμφωνούν μέχρι το τελευταίο
\tl{byte} για \tl{k} έως 4096, όπως αναφέρθηκε στην
ενότητα~\ref{sec:meth-traffic}. Η σύμπτωση δεν είναι τυχαία: εκτελείται το ίδιο
χρονοδιάγραμμα με ισοδύναμη στρατηγική συγχώνευσης, απλώς σε ρητά
διαχειριζόμενη μνήμη αντί για αυτόματη.

Για τη μεγαλύτερη τιμή του \tl{k} που εξετάζεται η συνθήκη καταρρέει. Τα
ταξινομημένα μπλοκ αποκτούν μέγεθος 131072 στοιχείων, δεκαέξι φορές μεγαλύτερο
από το πλακίδιο, οπότε το μεγαλύτερο μέρος των επιπέδων της πρώτης φάσης δεν
μπορεί πλέον να συγχωνευθεί και εκτελείται με περάσματα πάνω από την καθολική
μνήμη. Η μετρούμενη κίνηση αυξάνεται αντίστοιχα, από επταπλάσια σε
εικοσιπενταπλάσια έως πενηνταεξαπλάσια της αφελούς εκτίμησης ανάλογα με το
μέγεθος της εισόδου. Εδώ βρίσκεται η αιτία για την οποία ο συντελεστής
επανάληψης της ενότητας~\ref{sec:meth-traffic} δηλώνεται μόνο για \tl{k} πολύ
μικρότερο του \tl{n}, καθώς και ένα μέρος της εξήγησης για τη συμπεριφορά της
υλοποίησης στα μεγάλα \tl{k}.

\section{\tl{Map-reduce} με τοπικούς σωρούς}
\label{sec:gpu-mapreduce}

Ο αλγόριθμος είναι ο ίδιος με εκείνον της ενότητας~\ref{sec:cpu-mapreduce}: κάθε
νήμα διατηρεί σωρό μεγέθους \tl{k} με τα καλύτερα στοιχεία που έχει δει, η ρίζα
λειτουργεί ως κατώφλι, και τα επιμέρους αποτελέσματα συγχωνεύονται στο τέλος. Η
διαφορά βρίσκεται στην κλίμακα της αποσύνθεσης και στις συνέπειές της.

Στον επεξεργαστή τα νήματα είναι λίγα και κάθε ένα αναλαμβάνει συνεχόμενο τμήμα
της εισόδου. Στη μονάδα γραφικών τα νήματα είναι χιλιάδες και διατρέχουν την
είσοδο με βήμα ίσο με το συνολικό τους πλήθος, ώστε οι ταυτόχρονες προσβάσεις
γειτονικών νημάτων να αφορούν γειτονικές διευθύνσεις και να συγχωνεύονται σε
ενιαίες συναλλαγές μνήμης.

Κάθε νήμα διατηρεί δικό του σωρό, όχι κοινό ανά ομάδα. Οι σωροί δεν χωρούν στην
κοινή μνήμη, καθώς το γινόμενο του πλήθους των νημάτων επί \tl{k} υπερβαίνει
κατά πολύ τα διαθέσιμα 48 \tl{KB} ανά ομάδα, οπότε φυλάσσονται σε χώρο εργασίας
στην καθολική μνήμη. Αυτό είναι και το σημείο όπου η υλοποίηση αποκλίνει
ουσιωδώς από εκείνη του επεξεργαστή ως προς την κίνηση δεδομένων: ο χώρος
εργασίας είναι ανάλογος του γινομένου των νημάτων επί \tl{k}, και με χιλιάδες
νήματα ο όρος αυτός παύει να είναι αμελητέος.

Η συγχώνευση γίνεται σε δύο επίπεδα. Εντός κάθε ομάδας νημάτων, οι σωροί
συνενώνονται ανά ζεύγη σε δέντρο, με λογαριθμικό πλήθος βημάτων ως προς το
μέγεθος της ομάδας. Το αποτέλεσμα κάθε ομάδας συγχωνεύεται στη συνέχεια με τα
υπόλοιπα. Η δομή αυτή αποφεύγει τη σειριακή φάση που επιβαρύνει τον επεξεργαστή,
αλλά εισάγει το δικό της κόστος, καθώς κάθε βήμα του δέντρου διαβάζει και γράφει
σωρούς από την καθολική μνήμη.

\subsection{Όριο παραλληλίας για μεγάλα \tl{k}}
\label{sec:gpu-mapreduce-limit}

Η αποτελεσματικότητα του αλγορίθμου στηρίζεται εξ ολοκλήρου στην απόρριψη: ένα
στοιχείο που δεν ξεπερνά τη ρίζα του σωρού κοστίζει μία μόνο σύγκριση. Η ρίζα
όμως αποκτά νόημα ως κατώφλι μόνο αφού ο σωρός γεμίσει. Όσο περιέχει λιγότερα
από \tl{k} στοιχεία, κάθε νέο στοιχείο εισάγεται υποχρεωτικά και καμία απόρριψη
δεν συμβαίνει.

Από την παρατήρηση αυτή προκύπτει ένας ποσοτικός περιορισμός. Αν η είσοδος
μεγέθους \tl{n} μοιραστεί σε $P$ νήματα, κάθε νήμα βλέπει $n/P$ στοιχεία και ο
σωρός του γεμίζει μόνο αν $n/P \geq k$. Η συνθήκη αυτή ισοδυναμεί με
\[ P \leq \frac{n}{k} \]
και θέτει ανώτατο όριο στο πλήθος των νημάτων που μπορούν να συνεισφέρουν
ουσιαστικά. Πέραν αυτού, τα επιπλέον νήματα δεν επιταχύνουν τον υπολογισμό αλλά
τον επιβαρύνουν: γεμίζουν χώρο εργασίας με σωρούς που δεν φιλτράρουν ποτέ και
προσθέτουν επίπεδα στο δέντρο συγχώνευσης.

Ο περιορισμός είναι δεσμευτικός ακριβώς εκεί όπου η μονάδα γραφικών θα ήταν πιο
χρήσιμη. Για τη μεγαλύτερη τιμή του \tl{k} που εξετάζεται και για είσοδο
δεκαέξι εκατομμυρίων στοιχείων, ο λόγος $n/k$ είναι μόλις 128. Μια συσκευή με
χιλιάδες πυρήνες μπορεί επομένως να απασχολήσει ουσιαστικά εκατόν είκοσι οκτώ
νήματα, δηλαδή κλάσμα της τάξης του ενός τοις εκατό των διαθέσιμων. Η υποδομή
δεν υποαπασχολείται από αδυναμία υλοποίησης αλλά από τη δομή του ίδιου του
αλγορίθμου.

Η υλοποίηση αναγνωρίζει τη συνθήκη και αντιδρά σε αυτήν. Πριν από την εκκίνηση
υπολογίζεται αν οι σωροί προλαβαίνουν να γεμίσουν με πλήρη κατάληψη της
συσκευής. Όταν αυτό ισχύει, το πλήθος των ομάδων περιορίζεται σκόπιμα, αντίστροφα
ανάλογα του \tl{k}, ώστε κάθε νήμα να αναλάβει μεγαλύτερο μερίδιο εισόδου και ο
σωρός του να φιλτράρει αποδοτικά. Θυσιάζεται δηλαδή ονομαστική παραλληλία για
πραγματική απόρριψη. Όταν αντιθέτως οι σωροί δεν πρόκειται να γεμίσουν σε καμία
περίπτωση, η μείωση των ομάδων δεν θα προσέφερε τίποτα και η συσκευή
χρησιμοποιείται με πλήρη κατάληψη.

\section{Διαχείριση μνήμης και κίνηση δεδομένων}
\label{sec:gpu-memory}

\section{Υλοποίηση αναφοράς}
\label{sec:gpu-reference}

\subsection{Επιλογή με βάση την αναπαράσταση στη \tl{CUB}}
\label{sec:gpu-reference-cub}

\section{Περιοριστικοί παράγοντες}
\label{sec:gpu-bottlenecks}
